\chapter{稳定性与刚性方程}

\intro{
	不稳定性是数值分析的阿喀琉斯之踵。
	
	\rightline{——Germund Dahlquist}
}

\section{数值稳定性的基本概念}

\subsection{零稳定性与绝对稳定性}

数值稳定性讨论的是：局部截断误差和舍入误差在传播过程中是放大还是衰减。

\begin{mydefinition}{零稳定性}
	若存在 $h_0 > 0$ 使得对于任意 $h \in (0, h_0]$，方法的数值解对初始值的扰动是有界的（一致有界），则称该方法\textbf{零稳定}。
\end{mydefinition}

\begin{mydefinition}{绝对稳定性}
	对于试验方程 $y' = \lambda y$（$\operatorname{Re}(\lambda) < 0$），若一个数值方法在给定步长 $h$ 下产生的数值解 $y_n$ 满足 $\lim_{n\to\infty} y_n = 0$，则称方法对于 $z = \lambda h$ 是\textbf{绝对稳定}的。
\end{mydefinition}

\subsection{稳定函数与稳定域}

将单步法应用于试验方程 $y' = \lambda y$，可写为：
\[
y_{n+1} = R(z) \, y_n, \quad z = \lambda h
\]
其中 $R(z)$ 称为方法的\textbf{稳定函数}。

\begin{mydefinition}{稳定域}
	方法的绝对稳定域定义为：
	\[
	\mathcal{S} = \{ z \in \mathbb{C} \mid |R(z)| < 1 \}
	\]
\end{mydefinition}

\begin{theorem}{常用方法的稳定函数}
	\begin{enumerate}
		\item 显式 Euler：$R(z) = 1 + z$，稳定域为 $|1+z| < 1$（复平面上以 $-1$ 为心的单位圆）；
		\item 隐式 Euler：$R(z) = \dfrac{1}{1 - z}$，稳定域为 $|1-z| > 1$（包含整个左半平面）；
		\item 梯形法：$R(z) = \dfrac{1 + z/2}{1 - z/2}$，稳定域为整个左半平面；
		\item 经典 RK4：
		\[
		R(z) = 1 + z + \frac{z^2}{2} + \frac{z^3}{6} + \frac{z^4}{24}
		\]
	\end{enumerate}
\end{theorem}

\section{A-稳定性}

\begin{mydefinition}{A-稳定性（Dahlquist, 1963）}
	若方法的绝对稳定域包含整个左半复平面 $\{ z \in \mathbb{C} \mid \operatorname{Re}(z) < 0 \}$，则称该方法为\textbf{A-稳定}的。
\end{mydefinition}

\begin{theorem}{Dahlquist 第二障碍定理}
	\begin{enumerate}
		\item 显式线性多步法不可能是 A-稳定的。
		\item $A$-稳定的隐式线性多步法的最高阶数为 2。
		\item 2 阶 A-稳定的隐式线性多步法中，梯形法具有最小的误差常数。
	\end{enumerate}
\end{theorem}

\begin{remark}
	定理说明精确度与稳定性之间存在根本性矛盾：高阶方法无法做到完美（A-）稳定。这是数值分析中最深刻的结论之一。
\end{remark}

对于 Runge-Kutta 方法，不存在类似的阶数限制——存在任意高阶的 A-稳定 RK 方法（如 Gauss-Legendre 配点法）。

\section{刚性方程}

\subsection{刚性的概念}

\begin{mydefinition}{刚性方程（Stiff Equation）}
	若一个常微分方程的数值求解中，由于方程中同时包含快慢差异极大的时间尺度，导致显式方法为了保持稳定性而必须采用极小的步长（远小于精度所需的步长），则称该方程为\textbf{刚性}的。
\end{mydefinition}

\textbf{刚性的数学刻画：} 若 Jacobi 矩阵 $\partial \mathbf{f}/\partial \mathbf{y}$ 的特征值满足：
\[
\max_j |\operatorname{Re}(\lambda_j)| \gg \min_j |\operatorname{Re}(\lambda_j)|
\]
且所有 $\operatorname{Re}(\lambda_j) < 0$，则问题为刚性的。定义\textbf{刚性比}：
\[
S = \frac{\max_j |\operatorname{Re}(\lambda_j)|}{\min_j |\operatorname{Re}(\lambda_j)|}
\]
$S \gg 1$ 指示刚性问题。

\begin{example}{刚性方程示例}
	考虑系统：
	\[
	\begin{cases}
		y_1' = -1000 y_1 \\
		y_2' = -y_2
	\end{cases}
	\]
	特征值 $\lambda_1 = -1000$，$\lambda_2 = -1$。刚性比 $S = 1000$。用显式 Euler 法求解时，稳定条件要求 $h < 2/1000 = 0.002$，由精度要求看 $h < 2$ 就足够——计算效率损失 1000 倍！
\end{example}

\subsection{刚性方程的求解策略}

\begin{enumerate}
	\item \textbf{使用 A-稳定或 $A(\alpha)$-稳定的方法}：如 BDF、隐式 RK（Radau IIA）、梯形法等；
	\item \textbf{使用隐式格式}：在每一步需求解非线性代数方程组，通常用 Newton 迭代：
	\[
	\mathbf{y}_{n+1}^{(k+1)} = \mathbf{y}_{n+1}^{(k)} - \left[ I - h \gamma \, J(\mathbf{y}_{n+1}^{(k)}) \right]^{-1} \left[ \mathbf{y}_{n+1}^{(k)} - \mathbf{y}_n - h \gamma \, \mathbf{f}(t_{n+1}, \mathbf{y}_{n+1}^{(k)}) \right]
	\]
	其中 $J = \partial \mathbf{f}/\partial \mathbf{y}$ 为 Jacobi 矩阵；
	\item \textbf{自适应步长}：根据局部误差估计动态调整步长，在解变化剧烈处缩小步长，平缓处增大步长。
\end{enumerate}

\section{$A(\alpha)$-稳定性}

由于 A-稳定性对隐式多步法施加了严格的阶数限制，实际中常用更弱的定义：

\begin{mydefinition}{$A(\alpha)$-稳定性}
	若稳定域包含无限扇形区域 $\{ z \in \mathbb{C} \mid |\arg(-z)| \le \alpha, z \neq 0 \}$，则称方法为 \textbf{$A(\alpha)$-稳定}。
\end{mydefinition}

\begin{remark}
	$\alpha$ 越接近 $\pi/2$，稳定性越接近 A-稳定。BDF-3 的 $\alpha \approx 86^{\circ}$，BDF-4 的 $\alpha \approx 73^{\circ}$，BDF-5 的 $\alpha \approx 51^{\circ}$，BDF-6 不是 $A(\alpha)$-稳定的（$\alpha = 0$）。
\end{remark}

\section{非线性稳定性与强稳定性保持方法}

对于具有特殊性质的系统（如守恒系统、梯度流等），需要数值方法保持解的结构性质。例如：

\begin{enumerate}
	\item \textbf{辛几何积分器}：保持 Hamilton 系统的辛结构，适用于长时间仿真；
	\item \textbf{强稳定性保持（SSP）RK 方法}：保持总变差减小性质，适用于双曲守恒律的时间离散。
\end{enumerate}

\section{隐式 RK 方法与 Radau IIA}

对于刚性方程，隐式 Runge-Kutta（IRK）方法是一类强有力的工具。

\subsection{Radau IIA 方法}

Radau IIA 是具有优良稳定性性质的隐式 RK 族。其中 $s$ 级 Radau IIA 方法的稳定函数为 Padé 逼近 $(s-1, s)$：

\begin{myTheorem}{Radau IIA 方法的性质}
	$s$ 级 Radau IIA 方法具有以下性质：
	\begin{enumerate}
		\item 阶：$p = 2s - 1$；
		\item A-稳定且 L-稳定（$\lim_{z \to -\infty} R(z) = 0$）；
		\item 代数稳定（B-稳定）；
		\item 内部隐式求解的耦合使得每一步需解 $ns$ 维非线性系统。
	\end{enumerate}
\end{myTheorem}

对于刚性问题，3 级 Radau IIA 方法（5 阶 A-稳定）及其低阶版本被广泛采用，如 MATLAB ode15s 默认使用 BDF 方法，而 Julia 的 DifferentialEquations 包的 \texttt{RadauIIA5} 即基于此。

\section{步长控制与误差估计}

现代数值求解器中，步长控制是核心组件。通用策略为：

\begin{enumerate}
	\item \textbf{误差估计}：使用嵌入式对（如 RKF45, Dormand-Prince）或 Richardson 外推来估计局部误差；
	\item \textbf{步长调整}：若估计误差 $\Delta$ 超过容差 $\varepsilon$，拒绝该步并按公式缩小步长：
	\[
	h_{\text{new}} = \beta \cdot h \cdot \left( \frac{\varepsilon}{\Delta} \right)^{1/(p+1)}
	\]
	其中 $\beta \approx 0.8\sim0.9$ 为安全因子；
	\item \textbf{事件定位}：当系统中有不连续事件（如碰撞、开关），需精确确定事件发生的时刻。
\end{enumerate}

\begin{remark}
	在实际工程仿真中，选择正确的求解器和恰当的误差容差是平衡效率与精度的关键。一般经验：
	\begin{itemize}
		\item 非刚性：首选 \texttt{ode45}；
		\item 刚性：用 \texttt{ode15s}；
		\item 确定性问题且需要极高精度：\texttt{ode113}。
	\end{itemize}
\end{remark}
